\section{First-order rewriting, critical pairs and completion}

\begin{annexitem}{annex:p12.rewriting.closure}{thm:p12-rewrite-closure}{Closure under contexts, substitutions and repeated-variable replacement}
\begin{proof}
For a step \(s=D[l\tau]\to_R D[r\tau]=t\), applying a
substitution \(\sigma\) gives the same rule instance with substitution
\(\tau\sigma\), defined by
\(x(\tau\sigma)=(x\tau)\sigma\), in the context obtained by
applying \(\sigma\) to the variable leaves of \(D\).
Putting this entire step inside a further context is again an instance
of the definition of rewriting. Thus a step is preserved by substitution
and by an arbitrary outer context. Apply this observation to each step
of a finite reduction; a zero-step reduction remains a zero-step
reduction. Applying it to the two reductions to a common descendant
proves preservation of joinability.

For the last assertion, induct on the finite term \(u\). If \(u=x\),
use the stipulated reduction \(\sigma(x)\to_R^*\sigma'(x)\).
If \(u\) is another variable, its two instances are equal, so use
zero steps. If \(u=f(u_1,\ldots,u_k)\), the induction hypothesis
gives \(u_i\sigma\to_R^*u_i\sigma'\) for every child. Reduce the
first child inside the outer \(f\)-context, then the second, and
continue through the finite list. Already changed children are left
fixed. This gives \(u\sigma\to_R^*u\sigma'\). For a constant
\(k=0\), the argument is the zero-step case. Every occurrence of
\(x\) is thereby changed, even when there are several copies; if
there are none, all child reductions have length zero.
\end{proof}
\end{annexitem}

\begin{annexitem}{annex:p12.rewriting.newman}{thm:p12-newman}{The complete well-founded joining argument}
\begin{proof}
Write \(s>t\) when \(s\to_R^+t\), where the superscript denotes
one or more steps. This strict descendant relation is well-founded
under the termination hypothesis. Prove by well-founded induction on
\(s\) that any two finite reductions from \(s\) have joinable
endpoints.

Consider \(s\to_R^*t_1\) and \(s\to_R^*t_2\). If either
reduction has length zero, its endpoint is \(s\); the other endpoint
is a common descendant, using the other reduction and zero steps from
itself. Otherwise write their first steps as \(s\to_R u_1\) and
\(s\to_R u_2\). Local confluence supplies a term \(v\) with
\(u_1\to_R^*v\) and \(u_2\to_R^*v\).

Both \(u_1,u_2\) are strict descendants of \(s\). The induction
hypothesis at \(u_1\), applied to its reductions to \(t_1\) and
\(v\), supplies \(z_1\) with
\(t_1\to_R^*z_1\) and \(v\to_R^*z_1\). At \(u_2\) it
supplies \(z_2\) with \(t_2\to_R^*z_2\) and
\(v\to_R^*z_2\). Also \(v\) is a strict descendant of \(s\),
since \(s\to_R u_1\to_R^*v\) has at least one step. In
particular it cannot equal \(s\), which would give a cycle and
contradict termination. Apply the induction hypothesis at \(v\)
to obtain a common descendant \(z\) of \(z_1,z_2\).
Concatenation gives \(t_1\to_R^*z\) and \(t_2\to_R^*z\).
This proves the induction assertion for every \(s\), hence confluence.
\end{proof}
\end{annexitem}

\begin{annexitem}{annex:p12.rewriting.critical-pairs}{thm:p12-critical-pair}{All ordinary critical and variable overlap cases}
\begin{proof}
If \(R\) is locally confluent, the two components of each critical
pair are the endpoints of two one-step reductions from \(l_2\theta\).
Local confluence makes them joinable. Trivial pairs have the same
component twice and are joinable by zero steps.

Conversely, assume every critical pair is joinable. Consider an arbitrary
peak \(t\leftarrow_R s\to_R u\), with redex positions \(p,q\)
in \(s\). Each term and each left-hand side is finite, so the
position comparison has the following exhaustive cases.

\paragraph{Disjoint positions.}
If neither position is a prefix of the other, replacing the subterm at
\(p\) leaves the subterm at \(q\) unchanged, and conversely.
Perform on \(t\) the original \(q\)-step and on \(u\) the
original \(p\)-step. The resulting terms are literally equal: both
are obtained from \(s\) by making the same two replacements in
disjoint subtrees. This gives a one-step join on each side. No linearity
assumption is involved, since these are occurrences in the already
instantiated term.

\paragraph{A nonvariable overlap.}
Otherwise interchange the two sides if necessary so that \(p\) is
a prefix of \(q\). Factor out the context above \(p\). We may
therefore consider a root step \(l_2\sigma\to_R r_2\sigma\)
and an inner step at a relative position \(a\) of \(l_2\sigma\).
Take renamed copies of the two rules whose variable sets are disjoint
from one another and from the variables in the concrete peak. Suppose
first that \(a\) is a nonvariable position of the pattern \(l_2\).
For the inner rule \(l_1\to r_1\), its match \(\tau\) satisfies
\[
 (l_2|_a)\sigma=l_1\tau.
\]
Because the rule-variable sets are disjoint, their matching assignments
combine into one substitution \(\rho\) on their union. This
substitution unifies \(l_2|_a\) and \(l_1\). The first-order MGU
factorization theorem gives an MGU \(\theta\) and a substitution
\(\delta\) with \(x\rho=(x\theta)\delta\) on every variable
of the two left-hand sides. Variables present only outside the unified
subterms can be left fixed by the MGU and assigned by \(\delta\).
The condition \(\operatorname{Var}(r_i)\subseteq
\operatorname{Var}(l_i)\) extends the same identity to both right-hand
sides.

The two concrete reducts under the removed context are therefore
\[
 (r_2\theta)\delta
 \quad\hbox{and}\quad
 ((l_2[r_1]_a)\theta)\delta.
\]
The uninstantiated pair is a critical pair by definition. Its assumed
joining reductions remain valid after applying \(\delta\) and
restoring the outer context, by Lemma~\ref{thm:p12-rewrite-closure}.
They join the original endpoints. This includes \(a=\epsilon\),
different rules at the same root, and two renamed copies of one rule.

\paragraph{A reduction inside a substituted variable.}
The remaining nested case has the inner redex inside the instance of
a variable \(x\) occurring in \(l_2\). Write its position as
\(a=b c\), where \(l_2|_b=x\) and \(c\) is a position of
\(\sigma(x)\), possibly the root. The inner step changes
\(\sigma(x)\) to a term \(d\). Define \(\sigma'\) to agree
with \(\sigma\) on all variables except that \(\sigma'(x)=d\).
The inner-first branch has changed one occurrence of \(\sigma(x)\)
in \(l_2\sigma\). Reduce each of its other occurrences corresponding
to a leaf \(x\) of \(l_2\) by that same inner step. These leaves
have pairwise disjoint positions in the pattern, so changing one does
not remove or alter the others. After finitely many such steps the
result is exactly \(l_2\sigma'\), even if \(l_2\) is nonlinear.
The root rule now applies again and gives \(r_2\sigma'\).

The outer-first branch was \(r_2\sigma\). By the last assertion of
Lemma~\ref{thm:p12-rewrite-closure}, it reduces to
\(r_2\sigma'\) by reducing the copy of \(\sigma(x)\) at every
occurrence of \(x\) in \(r_2\). If the rule duplicates \(x\),
there are several copies and all are reduced. If it erases \(x\),
the two right-hand instances are already equal and this branch takes
zero further steps. Reinstating the outer context yields a join of
the original peak.

These are all cases: on following a position down the instantiated
left-hand side, either its endpoint is a function-symbol occurrence of
the pattern or the path first enters a variable instance. The proof
uses the ability to reduce every needed copy, and uses the root rule
at the altered substitution without checking an additional condition.
Those are precisely the unrestricted and unconditional hypotheses.
All one-step peaks are joinable, proving local confluence.
\end{proof}
\end{annexitem}

\begin{annexitem}{annex:p12.rewriting.normal-forms}{thm:p12-rewrite-normal-forms}{Decreasing reductions, normal forms and conversions}
\begin{proof}
If \(l\succ r\) for every rule, closure of the ordering under
substitution gives \(l\sigma\succ r\sigma\), and closure under
contexts gives \(C[l\sigma]\succ C[r\sigma]\). Thus every
rewrite step decreases the well-founded order. A nonempty set lacking
a rewrite-minimal member would also lack a \(\succ\)-minimal
member, so the rewrite relation is well-founded and terminating.

Well-founded induction proves that each term has a normal-form
descendant. A term with no successor is itself such a descendant.
Otherwise take one successor; by induction it has a normal-form
descendant, which is also a descendant of the original term. This is
an existential argument for each fixed term. With local confluence,
Newman's lemma gives confluence. If two normal forms descend from the
same term, confluence joins them; their irreducibility forces the joining
reductions to have length zero, so they are identical. Denote this
unique normal form by \(\operatorname{nf}_R(t)\).

For a step \(s\to_R t\), the normal forms of \(s\) and \(t\)
are normal descendants of \(s\), hence equal. The same equality
holds across a backward step. Induction on a finite conversion gives
equality of the endpoint normal forms. Conversely, if both terms reduce
to the same normal form, following the first reduction and the reverse
of the second gives a finite conversion between them.
\end{proof}
\end{annexitem}

\begin{annexitem}{annex:p12.rewriting.invariant}{thm:p12-completion-invariant}{Every registered equational-invariant case}
\begin{proof}
It is enough in each transformation to show that every old generating
equation follows from the new generators and every new generating
equation follows from the old. Taking the least equivalence closed
under contexts and substitutions then gives equality of the generated
relations.

\paragraph{Orient.}
Replacing the equation \(s\approx t\) by the rule \(s\to t\)
does not change its underlying equational generator. The order and
well-formed-rule side conditions govern the direction of rewriting,
not the symmetric generated equality.

\paragraph{Simplify.}
If \(s\to_R s'\), the retained rules already imply
\(s\equiv_R s'\). The old equation \(s\approx t\) therefore
implies the new \(s'\approx t\), and the new implies the old by
symmetry and transitivity with that same retained equality. The
argument for simplifying the other equation side is identical with
the side names interchanged.

\paragraph{Compose.}
Let \(R_0\) be the retained rules and let the distinguished rule be
\(s\to t\). The premise \(t\to_{R_0}t'\) gives
\(t\equiv_{R_0}t'\). Thus adjoining \(s\approx t\) or
\(s\approx t'\) to \(E\cup R_0\) generates the same relation:
each equation follows from the other and that retained equality.

\paragraph{Collapse.}
The distinguished rule \(s\to t\) is removed, and the remaining
rules \(R_0\) supply \(s\to_{R_0}s'\). Hence
\(s\equiv_{R_0}s'\) remains available after removal. In the old
theory it combines with \(s\approx t\) to yield \(s'\approx t\).
In the new theory the added equation \(s'\approx t\) combines
with \(s\equiv_{R_0}s'\) to recover \(s\approx t\).
Thus the removed rule is still an equational consequence, and the new
equation was already an old consequence. The argument covers every
Collapse instance permitted by the additional redex-eligibility
condition, since it only uses the retained-rule reduction.

\paragraph{Deduce.}
For a critical peak \(s\leftarrow_R u\to_R t\), the existing
rules imply \(s\equiv_R u\equiv_R t\). Its equation is therefore
already in the generated congruence, whether or not its two sides
are currently joinable by forward reductions. Adjoining the equation
preserves that congruence.

Every unchanged generator remains present in all five cases. Initially
the generated relation is \(\equiv_{E_0}\), since the rule set is
empty. Induction on the finite transformation sequence proves the
claimed invariant at every stage.
\end{proof}
\end{annexitem}

\begin{annexitem}{annex:p12.rewriting.completion}{thm:p12-completion-correctness}{The final equational decision theorem}
\begin{proof}
The invariant and empty final equation set give
\(\equiv_{E_0}=\equiv_R\). This latter relation equals the finite
conversion relation \(\leftrightarrow_R^*\). Indeed, every rewrite
step is a contexted substitution instance of an underlying equation,
so its two directions and every finite conversion belong to
\(\equiv_R\). Conversely, conversion is an equivalence relation;
Lemma~\ref{thm:p12-rewrite-closure}, applied step by step in either
direction, shows it is closed under contexts and substitutions. It
contains every generating rule equation, so the least such relation
\(\equiv_R\) is contained in conversion.

The fixed reduction ordering proves termination. Joinability of all
critical pairs proves local confluence by
Theorem~\ref{thm:p12-critical-pair}, including nonlinear and erasing
rules. Newman's lemma proves confluence. Thus
Lemma~\ref{thm:p12-rewrite-normal-forms} gives
\[
 s\equiv_{E_0}t
 \ \Longleftrightarrow\ s\mathrel{\leftrightarrow_R^*}t
 \ \Longleftrightarrow\
 \operatorname{nf}_R(s)=\operatorname{nf}_R(t).
\]
These deductions use the final conditions themselves, so they apply
in particular to a successful fair run without needing a further
claim that fairness alone ensures those conditions or termination.

If the final rule set is finite and terms and symbols are effectively
represented, one can find a redex by checking the finitely many positions
of the current finite term against the finite rule list. First-order
matching is decidable by reading the pattern: the first occurrence of
a variable records its target subterm and each repeated occurrence
must match the same subterm; symbol and arity mismatches reject.
Use a fixed order of positions and rules to select a step whenever one
exists. Termination forces this procedure to stop at a normal form.
Compute the two normal forms and compare them syntactically. The proved
equivalence makes this a decision of the original word problem.
Neither this algorithm nor the equivalence supplies an algorithm that
forces completion to terminate on every initial presentation.
\end{proof}
\end{annexitem}

\begin{annexitem}{annex:p12.rewriting.example}{thm:p12-completion-worked}{Ordering, critical pairs and normal forms in the worked completion}
\begin{proof}
Interpret the signature in \(\N\) by
\(a=2\), \(c=1\), \(b=0\) and \(f(n)=n+1\).
For a variable assignment \(\rho:V\to\N\), let
\(\llbracket t\rrbracket_\rho\) be recursive evaluation of a term.
Define \(s\succ t\) when
\(\llbracket s\rrbracket_\rho>
\llbracket t\rrbracket_\rho\) for every \(\rho\).
It is a strict transitive order. Any nonempty set of terms has a
member minimizing evaluation at the all-zero assignment, and that
member cannot have a \(\succ\)-smaller member in the same set.
Thus the ordering is well-founded.

Substitution stability follows by composing assignments: evaluation
of \(t\sigma\) at \(\rho\) is evaluation of \(t\) under
the assignment sending \(x\) to
\(\llbracket x\sigma\rrbracket_\rho\). A one-hole context in
this unary signature adds a fixed nonnegative integer to the hole's
value, so it preserves strict inequality. Hence \(\succ\) is a
reduction ordering. It orients the three ground rules because
\(2>1\), \(3>0\), and \(2>0\), respectively.

The two initial Orient steps produce \(a\to c\) and
\(f(a)\to b\). The peak at \(f(a)\) has endpoints \(b\)
and \(f(c)\), both irreducible at that stage and syntactically
distinct. Deduce therefore adds \(f(c)\approx b\), and the
ordering permits Orient to turn it into the third rule. The final
equation set is empty.

To enumerate all final critical pairs, a left-hand side is one of
\(a,f(a),f(c)\). At the root two of these unify only if they are
the same term, giving a trivial self-pair. The only other nonvariable
positions are the argument constants of \(f(a)\) and \(f(c)\).
Only the former unifies with a left-hand side, namely \(a\).
Its pair is exactly \((b,f(c))\), now joined by \(f(c)\to b\).
Thus all critical pairs are joinable. The completion theorem applies.

The reductions \(f(a)\to b\) and
\(f(f(a))\to f(b)\) give the two asserted equalities.
Neither \(b\) nor \(f(b)\) contains a redex, and they are distinct
terms. Their normal forms therefore differ, proving
\(f(b)\not\equiv_{E_0}b\). Finally, simplifying the initial
\(f(a)\approx b\) with \(a\to c\) before orienting it gives
the two-rule final set \(a\to c,f(c)\to b\); the same ordering
orients it and it has only trivial critical pairs. The invariant
establishes its identical original equational theory.
\end{proof}
\end{annexitem}
